So far we have drawn our pictures on paper ruled in small squares.
The \robot\ has the additional ability to change the size of the squares.
We can therefore enlarge the drawings by increasing the size of the squares before starting the drawing. It is also possible
to change the size of the squares while the \robot\ is drawing. This makes it possible to
draw segments of different lengths and arcs of circles whose radius can vary.
By choosing rectangular cells we can draw ellipses and choose
the angles of movement of the \robot\ freely.

\section{Square grid}
The standard grid shown in figure \ref{fig:grille} (page \pageref{fig:grille})
has a size of 12 pixels. We can change the
length of the side of a square of the grid with the command {\idt gs} (\underline{g}rid \underline{s}ize), which
takes as parameter the new length of the grid expressed in pixels. For example
{\idt gs(6)} sets a grid half the size of the usual grid whereas
{\idt gs(24)} doubles the size of the grid.

The length of the grid can also
have a decimal part, as in {\idt gs(7.5)} (seven and a half pixels). Such a
number, which is called a {\em floating-point number}\index{Floating-point number}, is written with a point
(and not with a comma as in French):
{\idt 6.0} (or {\idt 6.}) is the floating-point number six
whereas {\idt 6} is the integer six. The commands {\idt gs}, {\idt gsX} and {\idt gsY}
accept either an integer or a floating-point number: {\idt gs(6)} and {\idt gs(6.0)}
have the same effect\footnote{\ This is not the case in all languages: in \ocaml\ for example,
one must write {\idt gs\ 6.} with a floating-point number and convert integers
into floating-point numbers.}.
When the length of the grid is given by an integer expression such as
$(6\ \ast\ ${\idt i}$)$, we therefore simply write {\idt gs(6\ $\ast$\ i)}. We will study
floating-point numbers in detail in the chapter devoted to them.

\begin{exemple}[Superimposed squares]{\rm
 The program below draws superimposed squares, the \robot\ starting
 and ending at the lower left corner, facing north, the superimposition being
 obtained by choosing successive grid sizes of 6, 12, 18, \ldots, 84:

% \includePostScript{CarresSuperposes2}{0.95}{0.95}
% incorpore dans~:
 \input python/OverlappingSquares2.tex
\noindent $\Box$ }\end{exemple}


\section{Rectangular grid}

The \robot\ can walk on a rectangular grid whose horizontal side has a length
different from that of the
vertical side. This allows it to draw rec\-tan\-gles and el\-lip\-ses like these:
%de la figure \ref{fig:RecEl} (page \pageref{fig:RecEl}).
\begin{figure}[htbp]
 \includevspacePostScript[2.2cm]{RectangleEtEllipse40x20}{0.95}{0.95}
 \caption{Rectangle and ellipse on a $40 \times 20$ grid}
% \label{fig:RecEl}
\end{figure}

 The initial grid has a size of $12 \times 12$ counted in pixels. The command {\idt gsX(40)} gives
 a length of 40 pixels to the horizontal side of the grid without changing the length of the vertical side. After
 this command the grid therefore has a size of $40 \times 12$ pixels. In this notation $40 \times 12$ we give
 first the length of the horizontal side (called $X$) then the length of the vertical side (called $Y$)
 of the grid. The command {\idt gsY(20)} gives
 a length of 20 pixels to the vertical side of the grid without changing the length of the horizontal side. After
 this command the grid therefore has a size of $40 \times 20$ pixels. So we can change independently the
  horizontal size (with {\idt gsX}) or the vertical size (with {\idt gsY}) of the grid.
  We can return to a square grid by using the command {\idt gs}.
  For example the command {\idt gs(12)}
  returns to a standard $12 \times 12$ square grid. The rectangle and the ellipse
 above
  %de la figure \ref{fig:RecEl} (page \pageref{fig:RecEl})
  are built by the following program:\hfill
 \input python/RectangleAndEllipse.tex

 The figure
 below
%\ref{fig:vdavvg3060} (page \pageref{fig:vdavvg3060})
illustrates the effect of the commands
 {\idt cr()}, {\idt fd()} and {\idt cl()} with a rectangular grid of dimensions $30 \times 60$.
\begin{figure}[htbp]
 \includevspacePostScript[4.3cm]{VD+AV+VG-30x60}{0.95}{0.95}
 \caption{Curve to the right, move forward and curve to the left with a $30 \times 60$ grid}
% \label{fig:vdavvg3060}
\end{figure}
\begin{figure}[htbp]
 \includevspacePostScript[2.9cm]{Parallel+Ellipses}{0.95}{0.95}
 \caption{Parallelograms and ellipses}
 \label{fig:parel}
\end{figure}
\begin{exemple}{\rm
\renewcommand{\carrenoir}{\hspace{-4mm}\rule{1.5mm}{1.5mm}\hspace{2.5mm}}
The various segments and arcs of el\-lip\-ses drawn by the commands {\idt fd()}, {\idt cl()}
and {\idt cr()} are illustrated by the two programs below, which draw parallelograms and
ellipses by turning them a quarter turn around the centre of the grid, as can be seen in
fig\-ure \ref{fig:parel} (page \pageref{fig:parel}):\hfill
\input python/Parallelograms.tex
\input python/Ellipses.tex
\renewcommand{\carrenoir}{}
\noindent $\Box$}\end{exemple}

\needspace{6\baselineskip}\noindent \rule{137mm}{0.1mm}\par\nopagebreak %---------------------------------------------
\begin{exercices}{Use the possibilities offered by changing the size of the square grid
to draw:

\tiret the spirals and concentric circles of figure \ref{fig:spirales} (page \pageref{fig:spirales});

\tiret the frieze of figure \ref{fig:dfrisec} (page \pageref{fig:dfrisec});

\tiret the propeller and the ball of figure \ref{fig:heba} (page \pageref{fig:heba});

\tiret the snail of figure \ref{fig:escar} (page \pageref{fig:escar});

\tiret the stegosaurus and its baby of figure \ref{fig:steg2} (page \pageref{fig:steg2});

\tiret the family of dogs of figure \ref{fig:chiens} (page \pageref{fig:chiens});

\tiret the tiling of figure \ref{fig:PVdV} (page \pageref{fig:PVdV}) inspired by a motif by
van de Vecht (\cite{escher1}, page 26) shown in figure \ref{fig:mVdV} (page \pageref{fig:mVdV}).
\noindent $\Box$}\end{exercices}
\begin{exercices}{Use the possibilities offered by rectangular grids
to draw:

\tiret the frieze of overlapping diamonds of figure \ref{fig:frlos} (page \pageref{fig:frlos});

\tiret the bicycle of figure \ref{fig:velo} (page \pageref{fig:velo});

\tiret the eight-point compass rose
\ref{fig:rose-des-vents-8d} (page \pageref{fig:rose-des-vents-8d});

\tiret the tilings of figures \ref{fig:PVdV1l} (page \pageref{fig:PVdV1l}),
\ref{fig:PVdV2l} (page \pageref{fig:PVdV2l}) and \ref{fig:PVdV4l} (page \pageref{fig:PVdV4l}) obtained
from the tiling of figure \ref{fig:PVdV} (page \pageref{fig:PVdV}) by pro\-gres\-sive\-ly decreasing the size of the grid on the right-hand side
of the drawing, on the left-hand side and, in the last case, on all four sides of the rectangle, in the manner of the Dutch
engraver M. C. Escher (\cite{escher1}, \cite{escher2}).

\noindent $\Box$}\end{exercices}
\begin{figure}[htbp]
 \includevspacePostScript[2.8cm]{Spirales}{0.95}{0.95}
 \caption{Spirals and concentric circles}
 \label{fig:spirales}
\end{figure}
\begin{figure}[htbp]
 \includevspacePostScript[0.6cm]{DoubleFriseCrantee}{0.95}{0.95}
 \caption{Double notched frieze}
 \label{fig:dfrisec}
\end{figure}
\begin{figure}[htbp]
 \includevspacePostScript[4.3cm]{HeliceBalle}{0.95}{0.95}
 \caption{Propeller and ball}
 \label{fig:heba}
\end{figure}
\begin{figure}[htbp]
 \includevspacePostScript[4.8cm]{Escargot}{0.95}{0.95}
 \caption{Snail}
 \label{fig:escar}
\end{figure}
\begin{figure}[htbp]
 \includevspacePostScript[5.8cm]{StegosaureEtBebe}{0.95}{0.95}
 \caption{Stegosaurus and its baby}
 \label{fig:steg2}
\end{figure}
\begin{figure}[htbp]
 \includevspacePostScript[4.8cm]{ChiensEtChiots}{0.95}{0.95}
 \caption{Family of dogs out for a walk}
 \label{fig:chiens}
\end{figure}
\begin{figure}[htbp]
 \includevspacePostScript[2.5cm]{MotifVanDeVecht}{0.95}{0.95}
 \caption{Ornamental motif by N. J. van de Vecht}
 \label{fig:mVdV}
\end{figure}
\begin{figure}[htbp]
 \includevspacePostScript[2.2cm]{FriseLosangesImbriques}{0.95}{0.95}
 \caption{Frieze of overlapping diamonds}
 \label{fig:frlos}
\end{figure}
\begin{figure}[htbp]
 \includevspacePostScript[4.1cm]{velo}{0.95}{0.95}
 \caption{Bicycle}
 \label{fig:velo}
\end{figure}
\begin{figure}[htbp]
 \includevspacePostScript[14.6cm]{PavageVdVRegulier}{0.95}{0.95}
 \caption{Regular tiling inspired by an ornamental motif\label{fig:PVdV}}
 \centering{by N. J. van de Vecht}
\end{figure}
\begin{figure}[htbp]
 \includevspacePostScript[14.6cm]{PavageVdV1limite}{0.95}{0.95}
 \caption{Rectangular tiling with one limit}
 \label{fig:PVdV1l}
\end{figure}
\begin{figure}[htbp]
 \includevspacePostScript[14.6cm]{PavageVdV2limites}{0.95}{0.95}
 \caption{Rectangular tiling with two limits}
 \label{fig:PVdV2l}
\end{figure}
\begin{figure}[htbp]
 \includevspacePostScript[17.0cm]{PavageVdV4limites}{0.95}{0.95}
 \caption{Rectangular tiling with four limits}
 \label{fig:PVdV4l}
\end{figure}

\clearpage
%%%%%%%%%%%%%%%%%%%%%%%%%%%%%%%%%%%%%%%%%%%%%%%%%%%
\begin{corrige}{
%\label{cor:Frise-1}
\renewcommand{\carrenoir}{\hspace{-4mm}\rule{1.5mm}{1.5mm}\hspace{2.5mm}}
We give a few programs; the other programs are in the folder
{\idt book/\allowbreak 10\_grid}:
\input python/GeometricSpiral.tex
\input python/Snail2.tex
\renewcommand{\carrenoir}{}
\noindent $\Box$}\end{corrige}
\begin{corrige}{
We give the most difficult program, namely the tiling of figure
\ref{fig:PVdV4l} (page \pageref{fig:PVdV4l}) obtained
from the tiling of figure \ref{fig:PVdV} (page \pageref{fig:PVdV}) by pro\-gres\-sive\-ly decreasing the size of the grid on all
four sides of the rectangle. The other programs are
in the folder {\idt book/\allowbreak 10\_grid}:
\input python/VdVTiling4Limits.tex
\noindent $\Box$}\end{corrige}

%%%%%%%%%%%%%%%%%%%%%%%%%%%%%%%%%%%%%%%%%%%%%%%%%%%%%%%%%%%%%%%%%%%%%%%%%%%
\noindent \framebox[\textwidth]{
\begin{minipage}{\myparboxwidth}
\setlength{\parindent}{2em}
\small \sf
\vspace{\abovedisplayskip}
\section*{Addenda}
If all the grids of stopping points of the \robot\ were aligned on the centre of the screen, it could happen
that after a change of grid the \robot\ would no longer be on a point of the grid. This is why the
 grid is always relative to the \robot\ and not to the screen. When the size of the grid changes, all the points of the grid (which can be seen with
the command {\idt dg()}) move, except the one on which the \robot\ stands. Consequently no command
of the \robot\ (except one, namely {\idt gc()}) can make it leave one of the points of the current grid. For example the commands {\idt gl()}, {\idt gr()}, {\idt gt()} and
{\idt gb()} move the \robot\ horizontally or vertically to the point of the current grid (respectively to the left, to the right, at the top and
at the bottom) closest to the edge.
To be able to place the \robot\ anywhere on the screen precisely, one needs an
 absolute reference point. This absolute reference point is the centre of the screen. The command
 {\idt gc()} always brings the \robot\ back to the same point, located exactly at the centre of the screen.
 This command is therefore the only
 command of the \robot\ that is not performed on the same grid. If the centre
 of the screen is not on the current grid, the command {\idt gc()} recentres the grid exactly
 at the centre of the screen so as to avoid making the \robot\ leave its grid of movement.

To avoid possible misalignment problems, it is enough to program the drawing as if it were
on graph paper, choosing a minimal grid size (1 millimetre on graph paper)
so that all the changes of grid size are made when the \robot\ is at the centre and correspond
to multiples (5, 10, \ldots\ on graph paper) of the basic size. One can also proceed
empirically by experimenting, which is sometimes more fun.

\vspace{\abovedisplayskip}
\end{minipage}
}
%%%%%%%%%%%%%%%%%%%%%%%%%%%%%%%%%%%%%%%%%%%%%%%%%%%%%%%%%%%%%%%%%%%%%%%%%%%%
